\documentclass[a4paper,10pt]{article}

\usepackage[T1]{fontenc}
\usepackage[utf8]{inputenc}
\usepackage{lmodern}
\usepackage[margin=1.35cm,bottom=1.15cm]{geometry}
\usepackage[hidelinks]{hyperref}
\usepackage{tabularx}
\usepackage{array}
\usepackage{titlesec}
\usepackage{xcolor}
\usepackage{enumitem}

\pagestyle{empty}
\setlength{\parindent}{0pt}
\setlength{\parskip}{3pt}
\setlength{\tabcolsep}{0pt}

\titleformat{\section}{\large\bfseries}{}{0em}{}[\titlerule]
\titlespacing{\section}{0pt}{6pt}{4pt}

\newcolumntype{L}{>{\raggedright\arraybackslash}p{3.2cm}}
\newcolumntype{Y}{>{\raggedright\arraybackslash}X}
\newcolumntype{S}{>{\raggedright\arraybackslash}p{4.8cm}}
\newcolumntype{K}{>{\raggedright\arraybackslash}p{1.7cm}}
\newcolumntype{R}{>{\raggedright\arraybackslash}p{4.3cm}}

\setlist[itemize]{leftmargin=1.2em,itemsep=1pt,topsep=2pt,parsep=0pt,partopsep=0pt}

\begin{document}

\begin{center}
    {\LARGE \textbf{Kourosh Hasheminejad}}\\[4pt]
    {\large Computational Materials Scientist | Multiscale Modelling | Soft Matter}\\[6pt]
    Espoo, Finland \,|\,
    \href{mailto:kourosh.hasheminejad@aalto.fi}{kourosh.hasheminejad@aalto.fi}\\[2pt]
    \href{https://kourosh-h.github.io}{kourosh-h.github.io} \,|\,
    \href{https://github.com/kourosh-h}{GitHub} \,|\,
    \href{https://fi.linkedin.com/in/kourosh-hasheminejad-2577001aa/}{LinkedIn} \,|\,
    \href{https://scholar.google.com/citations?user=Ey7-XMsAAAAJ&hl=en}{Google Scholar}
\end{center}

\section{Research Profile}
Computational materials scientist with a D.Sc. from Aalto University, more than seven years of experience in molecular and mesoscale modelling. Expertise includes atomistic molecular dynamics, coarse-grained modelling, Python-based analysis, HPC workflow development, interfaces, nanocomposites, coatings, and structure--property relationships. Experienced in industry-linked research with Kemira and in connecting simulation-derived descriptors with experimental trends, and communication of computational results to researchers with different technical backgrounds. Interested in applying this expertise in customer-facing scientific collaborations and broader real-world materials applications.

\section{Education}
\begin{tabularx}{\linewidth}{@{}K Y@{}}
2026 & \textbf{Doctor of Science (D.Sc.)}, Chemistry and Materials Science, Aalto University, Finland \\
2019 & \textbf{Master of Science (M.Sc.)}, Materials Science and Engineering, K. N. Toosi University of Technology, Iran. Metallurgy-focused training including solidification and materials processing. \\
2016 & \textbf{Bachelor of Science (B.Sc.)}, Materials Science and Engineering, Shahid Beheshti University, Iran. Metallurgy-focused training in metallic materials, heat treatment, mechanical properties, casting, and coatings. \\
\end{tabularx}


\section{Professional Experience}

\textbf{Doctoral Researcher, Aalto University, Finland} \hfill \textbf{2021--2026}
\begin{itemize}
    \item Conducted multiscale molecular simulations, from atomistic molecular dynamics to mesoscale modelling, of thin films, coatings, nanocomposites, and interfacial systems.
    \item Developed structure--property analyses connecting local interactions and interfacial organization with morphology, transport, adhesion, and material performance.
    \item Carried out simulation tasks in industrially linked R\&D projects with \textbf{Kemira}, translating applied materials questions into computational models and connecting simulation trends with experimental observations.
    \item Developed automated Python workflows for simulation setup, trajectory analysis, quantitative post-processing, and reproducible analysis of large materials datasets.
    \item Built Bash- and Slurm-based workflows for running, chaining, and managing simulation campaigns on Aalto Triton and CSC HPC systems.
    \item Worked across modelling and experimental perspectives, communicating simulation assumptions, validation limits, and physically meaningful descriptors to multidisciplinary collaborators.
\end{itemize}

\textbf{Research Assistant, Aalto University, Finland} \hfill \textbf{2020--2021}
\begin{itemize}
    \item Prepared molecular simulation systems, established computational workflows, and carried out pilot studies for materials and interface modelling projects.
\end{itemize}


\section{Metallurgy \& Materials Engineering Foundation}
\begin{itemize}
    \item Laboratory training in tensile testing, fracture testing, steel heat treatment, and aluminum casting during B.Sc. studies.
    \item Coursework across B.Sc. and M.Sc. studies included solidification, metallic materials, heat treatment, mechanical behavior of materials, casting, and materials processing.
    \item B.Sc. project: electrochemical deposition of nickel--alumina composite coatings on steel.
    \item M.Sc. research: molecular modelling of polymer/graphene nanocomposites, including interface engineering and experimentally guided simulation of shape-memory behavior.
\end{itemize}

\section{Selected Research Contributions}
\begin{itemize}
    \item Developed experimentally guided atomistic simulations for PLA/graphene nanocomposites to study how interface engineering influences adhesion and shape-memory behavior.
    \item Developed atomistic and mesoscale models for polymer coatings and self-assembling systems, linking molecular-scale interactions to morphology and water transport.
    \item Built reusable analysis and HPC workflows for simulation campaigns, enabling systematic comparison across model parameters and material conditions.
\end{itemize}


\section{Technical Skills}
\begin{tabularx}{\linewidth}{@{}S Y@{}}

\textbf{Computational Materials Modelling} &
Classical molecular dynamics (LAMMPS, GROMACS), coarse-grained and mesoscale modelling (DPD), multiscale simulation, structure--property analysis \\[2pt]

\textbf{Metallurgy \& Structural Materials} &
Solidification and heat-treatment fundamentals, metallic materials and processing, mechanical-property testing fundamentals, fracture and tensile testing laboratory experience, aluminum casting, metallic coatings \\[2pt]

\textbf{Interfaces \& Materials Physics} &
Polymer nanocomposites, coatings, surfaces and interfaces, adhesion, hydration, molecular transport, morphology, self-assembly, interfacial structure \\[2pt]

\textbf{Data \& Automation} &
Python (NumPy, Pandas, scikit-learn, Matplotlib, MDAnalysis), Bash scripting, automated trajectory analysis, large-scale dataset post-processing \\[2pt]

\textbf{Scientific Computing} &
Linux/Unix, HPC environments (Slurm, Aalto Triton, CSC), Git/GitHub, Jupyter Notebooks, reproducible computational workflows \\[2pt]

\textbf{Visualization \& Analysis} &
OVITO, VMD, molecular trajectory analysis, density and morphology analysis, interface and transport descriptors \\[2pt]

\textbf{Emerging Methods} &
Machine learning for materials, basic quantum chemistry (ORCA, PSI4), initial experience with machine-learning interatomic potentials (MACE) \\

\end{tabularx}


\section{Teaching, Training \& Scientific Service}
\begin{tabularx}{\linewidth}{@{}R Y@{}}
\textbf{Supervision} & Advised 5 B.Sc. theses and undergraduate research projects at Aalto University. \\[2pt]
\textbf{Teaching} & Content Developer and Teaching Assistant for \textit{Soft Materials Modelling}, \textit{Computational Chemistry I}, and \textit{Personal Research Assignment} at Aalto University. \\[2pt]
\textbf{Training Materials} & Developed self-contained Jupyter Notebook exercises for teaching computational methods through clear instructions, reproducible environments, and ready-to-run workflows. \\[2pt]
\textbf{Community Building} & Founder and Lead of the \textit{Modeller's Council} (2022--2025), a bi-weekly forum for software demos, paper discussions, and simulation troubleshooting. \\[2pt]
\textbf{Peer Review} & Reviewer for the \textit{Journal of Colloid and Interface Science} (2024). \\[2pt]
\textbf{Pedagogical Training} & Aalto University pedagogical training: \textit{Teaching Assistant as a Learning Instructor} (2024). \\
\end{tabularx}

\section{Publications}
\textbf{7 peer-reviewed publications}, including \textbf{4 as first author}; \textbf{1 submitted first-author manuscript}; \textbf{148 citations}; \textbf{h-index 6}.
\vspace{0.3em}

\small
\begin{tabularx}{\linewidth}{@{}K Y@{}}
2026 & \textbf{Hasheminejad, K.}; Tolmachev, D.; Sammalkorpi, M. Hydration-enhanced chain length effects in short alkyl chain self-assembled monolayers on silica. Submitted to \textit{Journal of Colloid and Interface Science}. \\
2025 & \textbf{Hasheminejad, K.}; et al. Molecular level dependencies of water permeability in di-block copolymer assemblies. \textit{Computational Materials Science} \textbf{258}, 114046. \\
2023 & \textbf{Hasheminejad, K.}; et al. Cracking polymer coatings of paper-like surfaces: Control via block co-polymer structure and system composition. \textit{Applied Surface Science} \textbf{640}, 158324. \\
2023 & Scacchi, A.; \textbf{Hasheminejad, K.}; et al. Controlling self-assembling co-polymer coatings via co-polymer block length ratio. \textit{Journal of Colloid and Interface Science} \textbf{640}, 809--819. \\
2022 & Amini, M.; \textbf{Hasheminejad, K.}; et al. Engineering the shape memory parameters of graphene/polymer nanocomposites. \textit{Smart Materials and Structures} \textbf{31}, 025010. \\
2021 & \textbf{Hasheminejad, K.}; et al. Tailoring adhesion characteristics of PLA/graphene nanocomposites by end-grafted polymer chains. \textit{European Polymer Journal} \textbf{148}, 110351. \\
2020 & Amini, M.; \textbf{Hasheminejad, K.}; et al. Experimentally guided MD simulation to enhance the shape memory behavior of polymer-based nanocomposites. \textit{Composites Part A} \textbf{138}, 106055. \\
2020 & \textbf{Hasheminejad, K.}; et al. Enhanced interfacial characteristics in PLA/graphene composites through numerically-designed interface treatment. \textit{Applied Surface Science} \textbf{502}, 144150. \\
\end{tabularx}
\normalsize

\section{Conference \& Workshop Presentations}
Participated in \textbf{6} conferences and workshops, delivering \textbf{4 oral presentations} and \textbf{2 posters} on computational materials science, soft matter, and polymer systems.
\vspace{0.3em}

\small
\begin{tabularx}{\linewidth}{@{}K Y@{}}
2024 & \textbf{European Colloid and Interface Society (ECIS)}, Copenhagen, Denmark. \textbf{Oral presentation:} \textit{Self-assembling layered polymeric coatings: A dissipative particle dynamics study}. \\
2024 & \textbf{Nordic Polymer Days}, Helsinki, Finland. \textbf{Oral presentation:} \textit{Dissipative particle dynamics simulations of self-assembling layered polymer coatings}. \\
2024 & \textbf{Computational Chemistry Days}, Jyväskylä, Finland. \textbf{Oral presentation:} \textit{A computational take on designing self-assembling block copolymer films}. \\
2024 & \textbf{Chemistry Days Young Researcher Symposium (KP2024-YRS)}, Helsinki, Finland. \textbf{Poster presentation:} \textit{Designing self-assembling polymeric films by dissipative particle dynamics simulations}. \\
2023 & \textbf{Challenges in the Physics of Active and Biological Matter Workshop}, Espoo, Finland. \textbf{Poster presentation:} \textit{Mesoscale modelling of self-assembling polymeric coatings}. \\
2022 & \textbf{GRK2516 Workshop, Controlling Structure Formation in Soft Materials}, Mainz, Germany. \textbf{Oral presentation:} \textit{Computational design of self-assembling polymeric coatings}. \\
\end{tabularx}
\normalsize

\end{document}
